\documentclass{../../test-template}
% Якщо шаблон ще не підключає hyperref, а пункти виходять із зайвими відступами,
% розкоментуйте:
% \usepackage{hyperref}
% \usepackage{enumitem}

\newcommand{\tablelink}{\href{https://docs.google.com/spreadsheets/d/1JiSbx1sBKMNnKzc8RCL8gHOllW0FGwNVWB05o250l_k/edit?usp=sharing}{таблиці розподілу тем}}
\newcommand{\deadline}{\underline{\hspace{3cm}}} % термін здачі

\begin{document}

\begin{center}
    \textbf{\Huge{Практичне завдання}} \\[4pt]
    \large Дослідження сучасних протоколів доведення з нульовим розголошенням
\end{center}

\section*{Мета роботи та постановка задачі}

Мета роботи --- навчитися аналізувати, як протоколи доведення з нульовим
розголошенням (ZK) застосовуються в реальних системах: від постановки
прикладної задачі до вибору конкретної системи доведення, її адаптації
та оцінки безпеки.

Кожен студент обирає один проєкт або технологію з переліку варіантів
(див. останній розділ) і записує себе в таблицю за посиланням \tablelink.
Одну тему може взяти лише один студент\footnote{Тему можна запропонувати
самостійно за умови погодження з викладачем.}.

Результат роботи --- письмовий звіт і доповідь із презентацією у форматі
міні-лекції для одногрупників.

\section*{План дослідження}

Дослідження складається з трьох частин.

\begin{enumerate}
    \item \textbf{Опис проєкту.} Яку задачу розв'язує проєкт і навіщо йому
    доведення з нульовим розголошенням: що саме доводиться, хто є доказувачем
    і верифікатором, яка інформація приховується. Яке рішення
    використовувалося до цього (або використовується в аналогах) і чому
    від нього відмовилися (якщо таке було).

    \item \textbf{Опис протоколу.} Яка система доведення лежить в основі
    (арифметизація, схема зобов'язань, перетворення Фіата--Шаміра, потреба
    в довіреному налаштуванні тощо) та її ключові властивості. Чи модифіковано
    протокол порівняно з початковою конструкцією; які особливості та відмінності
    має його використання саме в цьому проєкті.

    \item \textbf{Аналіз.} Можливі недоліки протоколу: криптографічні
    припущення, відомі вразливості чи інциденти, обмеження продуктивності.
    Чи можна замінити протокол на інший і чого це коштуватиме.
    Якщо доступні бенчмарки (розмір доведення, час доведення та верифікації),
    наведіть їх із посиланням на джерело або спробуйте оцінити самостійно (за можливості).
\end{enumerate}

\section*{Оформлення результатів}

Оцінюються три складові: звіт, презентація та виступ.

\paragraph{Звіт} оформлюється у вільному стилі, але розділи мають відповідати
трьом частинам плану дослідження. Здається у форматі PDF.
\begin{itemize}
    \item шрифт Times New Roman, 12~pt, одинарний міжрядковий інтервал;
    \item поля: ліве 2,5~см, праве 1,5~см, верхнє та нижнє 2~см;
    \item обсяг --- від 3 до 8 сторінок (без титульної сторінки та списку джерел);
    \item обов'язковий список використаних джерел з посиланнями на першоджерела
    (статті, специфікації, документацію, звіти аудитів).
    \item обов'язково зазначається використання ШІ в дослідженні, вказуючи модель. Якщо оцінки були побудовані з допомогою ШІ, то необхідно вказати про це (однак краще так не робити).
\end{itemize}

\paragraph{Презентація} виконується у довільному стилі, зручному доповідачеві,
і розрахована на доповідь тривалістю 7--10 хвилин. Доповідь має бути
міні-лекцією, що охоплює основні результати дослідження.

\section*{Критерії оцінювання}

\begin{center}
\begin{tabular}{|l|c|}
    \hline
    \textbf{Складова} & \textbf{Максимум балів} \\ \hline
    Звіт         & 10 \\ \hline
    Презентація  & 5  \\ \hline
    Виступ       & 5  \\ \hline
    \textbf{Разом} & \textbf{20} \\ \hline
\end{tabular}
\end{center}

\noindent Звіт оцінюється за повнотою опису проєкту (2 бали), коректністю
та глибиною опису протоколу (4 бали), якістю аналізу (3 бали)
та оформленням і джерелами (1 бал).

\medskip
\noindent \textbf{Термін здачі:} \textbf{зазначений у класрумі}. За кожен повний тиждень
запізнення оцінка знижується на 1 бал, максимально на -4 бали. Робота виконана із запізненням більше ніж 4 тижні оцінюється у 0 балів.

\section*{Перелік варіантів}

\subsection*{Приватність та ідентичність}
\begin{enumerate}
    \item Zcash: Sapling (Groth16) та Orchard (Halo~2).
    \item Tornado Cash та Privacy Pools.
    \item Semaphore та World ID.
    \item Sui zkLogin: доведення над OpenID-токенами.
    \item ZK-паспорти та zk-email (Self, ZKPassport, Anon Aadhaar).
    \item Google Longfellow ZK: перевірка віку над цифровими документами (Ligero + sumcheck).
    \item Анонімні облікові дані: BBS-підписи та Signal zkgroup.
    \item Aztec Network: приватні смарт-контракти (Noir, UltraHonk).
\end{enumerate}

\subsection*{zkVM та rollup-рішення}
\begin{enumerate}
    \setcounter{enumi}{8}
    \item RISC Zero: STARK для RISC-V та обгортка в Groth16.
    \item SP1 Hypercube: zkVM на основі sumcheck.
    \item Starknet: Cairo та S-two (Circle STARK).
    \item Scroll / OpenVM (Plonky3).
    \item Mina Protocol: Kimchi та рекурсія Pickles.
\end{enumerate}

\subsection*{Folding-схеми та постквантові доведення}
\begin{enumerate}
    \setcounter{enumi}{13}
    \item LatticeFold.
    \item Nova та атака на коректність його реалізації (2023).
    \item Neo / SuperNeo: решітковий folding над малими полями.
    \item Агрегація постквантових підписів в Ethereum (leanSig, leanMultisig).
    \item LaBRADOR та Greyhound (бібліотека LaZer).
\end{enumerate}

\subsection*{Bitcoin, Monero та інші застосування}
\begin{enumerate}
    \setcounter{enumi}{18}
    \item Monero: Bulletproofs+ та FCMP++.
    \item BitVM2 / Citrea: перевірка SNARK-доведень у Bitcoin.
    \item Неінтерактивний DKG з ZK-доведеннями (Groth21, Internet Computer).
    \item zkTLS: TLSNotary, Reclaim.
    \item Proof of reserves криптобірж (Binance, Summa).
\end{enumerate}

\end{document}